%% kapa: exam=1 round=26 year=2015 period=1 form=A subject=civil_law subject_ko=민법 seq=16 number=16
%% source: Q-Net 감정평가사 기출문제 (cst00309) artlSeq=5138001 file=extracted/1008/5138001/제26회_감정평가사 1차_기출/제26회_감정평가사 1차_기출/1교시_감평_A형.pdf
%% conversion: kordoc 4.24.51 -> tools/gen1.py (auto); review=proofread-sonnet; flags=none
\begin{question}{16}
법률행위의 부관에 관한 설명으로 옳지 않은 것은? (다툼이 있으면 판례에 따름)
\begin{choices}
  \item 임대인이 생존하는 동안 임대하기로 하는 계약은 기한부 법률행위이다.
  \item 해제조건부 법률행위에서 그 조건이 이미 성취할 수 없는 것인 경우 그 법률행위는 무효로 한다.
  \item 조건의 성취로 인하여 불이익을 받을 당사자가 신의 성실에 반하여 조건의 성취를 방해한 경우, 조건이 성취된 것으로 의제되는 시점은 이러한 행위가 없었더라면 조건이 성취되었으리라고 추산되는 시점이다.
  \item 어느 법률행위에 어떤 조건이 붙어 있었는지 여부는 그 조건의 존재를 주장하는 자가 입증하여야 한다.
  \item 선량한 풍속 기타 사회질서에 위반한 조건이 붙은 법률행위는 무효로 한다.
\end{choices}
\end{question}
\answer{2}
